\sect{Приложение В. Forward-рекурсия для гибридной полу-марковской
модели}{appendix-hsmm}
%
В настоящем приложении приводится прямая рекурсия для гибридной
архитектуры из~\Secref{math}, в которой множество состояний
разделено на две непересекающиеся группы:
$\mathcal{S}_{\text{semi}}$~--- полу-марковские состояния с явной
плотностью длительности $d_j(u)$, и $\mathcal{S}_{\text{mark}}$~---
обычные марковские с геометрической длительностью, заданной
параметром выхода~$p_j$ согласно~\Eqref{cont-implicit-duration}.
Псевдокод следует логике, описанной у Конта~\cite[\S~IV]{Cont2010}.

\par\addvspace{2pt}
{\centering
\begin{minipage}{0.97\linewidth}
\footnotesize
\textbf{Алгоритм 3.} Forward для гибридной HHSMM
(\cite[уравн.~3--4]{Cont2010}). \\[1pt]
\rule{\linewidth}{0.4pt}\vspace{1pt}\\
\textbf{Вход:} $A$, эмиссии $b_j(\cdot)$, длительности
$d_j(\cdot)$, верхняя граница $U_{\max}$ \\
\textbf{Выход:} forward-переменные $\alpha_j(t)$ \\
1: $\alpha_j(1) \gets \pi_j\,b_j(o_1)$ для всех~$j$ \\
2: \textbf{для} $t = 2, 3, \dots$ \textbf{выполнить} \\
3: \quad \textbf{для каждого} $j \in \mathcal{S}_{\text{semi}}$
   \textbf{выполнить} \\
4: \quad\quad $\alpha_j(t) \gets \max\limits_{u\le \min(t-1, U_{\max})}\!
   \Bigl\{\sum_{i\ne j}\alpha_i(t-u)\,A_{ij}\,d_j(u)\,
   \prod_{\tau=t-u+1}^{t} b_j(o_\tau)\Bigr\}$ \\
5: \quad \textbf{конец для} \\
6: \quad \textbf{для каждого} $j \in \mathcal{S}_{\text{mark}}$
   \textbf{выполнить} \\
7: \quad\quad $\alpha_j(t) \gets b_j(o_t)\,\max_i [A_{ij}\,
   \alpha_i(t-1)]$ \\
8: \quad \textbf{конец для} \\
9: \quad \textbf{выдать} $\hat\varphi(t) \gets
   \argmaxop_j \alpha_j(t)$ \\
10: \textbf{конец для} \\
\rule{\linewidth}{0.4pt}
\end{minipage}\par}
\addvspace{4pt}

Сложность рекурсии для полу-марковской ветки
составляет $O(N^2\,U_{\max})$ на шаг против $O(N^2)$ для
обычной HMM; для марковской ветки сложность остаётся
прежней. На практике используется ограничение $U_{\max}$,
определяемое максимально допустимой длительностью ноты в
произведении (типично несколько секунд). Полное обоснование
рекурсии и связь с приложением \cite[Appendix~B]{Cont2010}
требуют отдельного анализа, в настоящей работе не приведённого.

\paragraph*{Замечание о форме плотностей.} Для практической
реализации существенен выбор параметрической формы
$d_j(u)$. Конт~\cite[\S~II-C]{Cont2010} использует две основные
формы: компаунд-распределение
\begin{equation}
P(U=u) = \sum_{n=1}^{r-1}\binom{u-1}{n-1}(1-p-q)^{u-n}q^{n-1}p
+ \binom{u-1}{r-1}(1-p-q)^{u-r}q^{r-1}(p+q),
\eqlab{compound-occupancy}
\end{equation}
и его частный случай при $p=0$~--- отрицательное
биномиальное:
\begin{equation}
P(U=u) = \binom{u-1}{r-1}q^r(1-q)^{u-r}.
\eqlab{negative-binomial}
\end{equation}
Параметр $r$ имеет музыкально осмысленную интерпретацию
числа «микро-состояний» внутри одной ноты, $p$ и $q$~---
вероятности выхода и продвижения внутри ноты соответственно.
Альтернативно используются гамма- или log-нормальные
распределения; выбор формы должен соответствовать наблюдаемому
эмпирическому распределению длительностей в обучающем
датасете.
